\documentclass[a4paper,11pt]{article}

\usepackage[T1]{fontenc}
\usepackage{lmodern}
\usepackage[a4paper,margin=2.7cm]{geometry}
\usepackage{amsmath,amssymb}
\usepackage{xcolor}
\usepackage{booktabs}
\usepackage{hyperref}

\usepackage[
  shortcuts,
  theorems,
  draft,
  todoextension=5cm,
  framedtheorems
]{tssuite}

\hypersetup{
  colorlinks=true,
  linkcolor=blue!60!black,
  urlcolor=blue!60!black
}

\title{\texttt{tssuite}: a small text utility package}
\author{Example document and informal guide}
\date{\today}

% ------------------------------------------------------------------------------
% Example definitions used below
% ------------------------------------------------------------------------------

\DefineScName[\Cpp]{C++}
\DefineScName[\LatexThree]{LaTeX3}
\DefineScName{GitHub}
\DefineScName{Visual Studio Code}
\DefineScNames{Python, Rust, LuaLaTeX}

\CreateTextField{\AuthorName}{author}
\CreateTextField{\ProjectTitle}{project}
\CreateTextField{\DocumentVersion}{version}

\AuthorName{Alice Example}
\ProjectTitle{tssuite demonstration}
\DocumentVersion{0.3}

\NewNumberedTheorem{theorem}{Theorem}
\NewNumberedTheorem{lemma}{Lemma}
\NewStarredTheorem{remark}{Remark}
\NewDocumentTheorem{definition}{Definition}

\MakeFramedTheorem{ftheorem}{theorem}
\MakeFramedTheorem{fdefinition}{definition}

\NewNumberedTheorem{sectiontheorem}{Section Theorem}[parent=section]
\NewNumberedTheorem{sharedlemma}{Shared Lemma}[sibling=theorem]
\NewNumberedTheorem{plainexample}{Plain Example}[style=definition]

% ------------------------------------------------------------------------------
% Local helpers for this document only
% ------------------------------------------------------------------------------

\newcommand{\pkg}[1]{\texttt{#1}}
\newcommand{\cs}[1]{\texttt{\textbackslash#1}}

\newenvironment{codeblock}
  {\begin{quote}\small\begin{verbatim}}
  {\end{verbatim}\end{quote}}

\begin{document}

\maketitle

This document is meant to be read as a simple introduction to the
\pkg{tssuite} package.  It is not a formal reference manual.  Instead, it shows
the usual commands in context and explains the intended workflow in ordinary
text.  The package collects a few small helpers that are useful in longer
documents: compact list styles, small-caps shortcuts, named text fields, todo
notes, theorem declarations, and optionally framed theorem environments.

To use the full package, load it like this:

\begin{verbatim}
\usepackage[
  shortcuts,
  draft,
  todoextension=5cm,
  theorems,
  framedtheorems
]{tssuite}
\end{verbatim}

Build this example from the \texttt{examples} directory with
\texttt{latexmk -r latexmkrc example.tex}.  The build configuration keeps
auxiliary files in \texttt{.build/} and writes the final PDF beside this source
file.

The options can also be used separately.  The \verb|shortcuts| option enables
commands such as \verb|\DefineScName| and \verb|\DefineScNames|.  The
\verb|draft| option loads todo-note support and extends both sides of the page
to accommodate margin notes; the per-side extension width is controlled by the
\verb|todoextension| option (default \texttt{5\,cm}).  The \verb|theorems|
option enables the theorem declaration helpers.  The \verb|framedtheorems|
option enables framed theorem wrappers and also turns on theorem support.

The package defines a few list styles for \pkg{enumitem}.  They are used as
ordinary list options.  The most common one is \verb|compact|, which removes
extra vertical spacing:

\begin{verbatim}
\begin{itemize}[compact]
  \item first item
  \item second item
\end{itemize}
\end{verbatim}

For example:

\begin{itemize}[compact]
  \item first item;
  \item second item;
  \item third item.
\end{itemize}

There is also \verb|widemargins|, useful when a list should have a little more
visual structure:

\begin{enumerate}[widemargins]
  \item This item has more generous indentation.
  \item This style can make longer lists easier to scan.
\end{enumerate}

Small-caps shortcuts are useful when the same names or technical terms appear
many times in a document.  Instead of writing \verb|\textsc{...}| repeatedly,
define the name once in the preamble.  When the printed text is not a good
command name, pass the command explicitly:

\begin{verbatim}
\DefineScName[\Cpp]{C++}
\DefineScName[\LatexThree]{LaTeX3}
\end{verbatim}

These declarations create \verb|\Cpp| and \verb|\LatexThree|.  In the document
they print as \Cpp{} and \LatexThree.  When the text itself can be used to form
a command name, omit the optional argument:

\begin{verbatim}
\DefineScName{GitHub}
\DefineScName{Visual Studio Code}
\end{verbatim}

This creates \verb|\GitHub| and \verb|\VisualStudioCode|, which print as
\GitHub{} and \VisualStudioCode.  The package preserves the capitalization of
the original text when creating the command name.  Several names can be declared
at once:

\begin{verbatim}
\DefineScNames{Python, Rust, LuaLaTeX}
\end{verbatim}

This creates \verb|\Python|, \verb|\Rust|, and \verb|\LuaLaTeX|, which print as
\Python, \Rust, and \LuaLaTeX.

Text fields provide a small mechanism for storing and reusing named pieces of
text.  A field is created with a setter command and a field name:

\begin{verbatim}
\CreateTextField{\AuthorName}{author}
\CreateTextField{\ProjectTitle}{project}
\CreateTextField{\DocumentVersion}{version}
\end{verbatim}

The first argument is the command that will later set the value.  The second
argument is the field name used when printing the value.  After declaring the
fields, assign values like this:

\begin{verbatim}
\AuthorName{Alice Example}
\ProjectTitle{tssuite demonstration}
\DocumentVersion{0.3}
\end{verbatim}

The values can then be printed anywhere with \verb|\PrintTextField|.  For
example, this document currently has author
\PrintTextField{author}, project title \PrintTextField{project}, and version
\PrintTextField{version}.  A field can be updated later simply by calling its
setter command again.  For instance, after

\begin{verbatim}
\DocumentVersion{0.4}
\end{verbatim}

the printed version changes from the previous value to the new one.

\DocumentVersion{0.4}

The current version is now \PrintTextField{version}.  This is useful for
metadata such as author names, project titles, version numbers, course names,
or report identifiers.

Todo notes are available in \verb|draft| mode.  The \verb|todoextension| key
adds room on each side of the page; the default is \texttt{5\,cm}.  They use
the usual
\pkg{todonotes} command:

\begin{verbatim}
\todo{Rewrite this paragraph later.}
\end{verbatim}

Here is a sample todo note:

\todo{Rewrite this paragraph later.}

The theorem helpers are designed to make declarations short and consistent.
A numbered theorem is declared with:

\begin{verbatim}
\NewNumberedTheorem{theorem}{Theorem}
\end{verbatim}

After that, the environment can be used normally:

\begin{theorem}
If $a=b$ and $b=c$, then $a=c$.
\end{theorem}

Another numbered environment can be created in the same way:

\begin{verbatim}
\NewNumberedTheorem{lemma}{Lemma}
\end{verbatim}

\begin{lemma}
If $x>0$ and $y>0$, then $x+y>0$.
\end{lemma}

For unnumbered theorem-like material, use:

\begin{verbatim}
\NewStarredTheorem{remark}{Remark}
\end{verbatim}

\begin{remark}
This is an unnumbered remark.  It is useful for comments, explanations, or
observations that should not participate in theorem numbering.
\end{remark}

The command \verb|\NewDocumentTheorem| creates a paired theorem environment:
one numbered form and one unnumbered starred form.  For example,

\begin{verbatim}
\NewDocumentTheorem{definition}{Definition}
\end{verbatim}

creates both \verb|definition| and \verb|definition*|.

\begin{definition}
A prime number is a positive integer greater than $1$ with exactly two positive
divisors.
\end{definition}

When \verb|framedtheorems| is enabled, existing theorem environments can be
wrapped in frames.  First define the ordinary theorem environment, then define
a framed wrapper:

\begin{verbatim}
\MakeFramedTheorem{ftheorem}{theorem}
\MakeFramedTheorem{fdefinition}{definition}
\end{verbatim}

The first argument is the new framed environment.  The second argument is the
already existing theorem-like environment that should be placed inside the
frame.

\begin{ftheorem}
For every integer $n$, if $n$ is even, then $n^2$ is even.
\end{ftheorem}

A framed theorem may also have a title:

\begin{ftheorem}[Transitivity]
If $x<y$ and $y<z$, then $x<z$.
\end{ftheorem}

The same idea works for definitions:

\begin{fdefinition}[Group]
A group is a set $G$ together with a binary operation satisfying associativity,
identity, and inverses.
\end{fdefinition}

The theorem commands in \pkg{tssuite} pass their optional key list directly to
\pkg{keytheorems}.  This means ordinary \pkg{keytheorems} options such as
\verb|parent=section|, \verb|sibling=theorem|, and \verb|style=definition| can be
used through the shorter wrapper commands.

For example, this theorem is numbered inside the current section:

\begin{sectiontheorem}[note=Section-local numbering,label=thm:section-local]
The number of this theorem includes the section number.
\end{sectiontheorem}

The next environment shares the counter of the ordinary theorem environment:

\begin{sharedlemma}[note=Shared counter]
This lemma uses the same counter sequence as \verb|theorem|.
\end{sharedlemma}

The next example uses the \verb|definition| style, so its body is upright rather
than italic.

\begin{plainexample}[note=Definition-style body]
This is useful for examples, definitions, and other explanatory material.
\end{plainexample}

The optional argument of a theorem can also use key-value syntax.  Here we set a
label and a note at the same time:

\begin{theorem}[label=thm:label-test,note=A labelled theorem]
This theorem can be referenced later.
\end{theorem}

For example, Theorem~\ref{thm:label-test} was labelled using a key in the
optional argument.

The \verb|store| key can be used to save a theorem and print it again later.  This
may require compiling the document more than once.

\begin{theorem}[store=stored-demo,note=Stored theorem]
This theorem is stored by \pkg{keytheorems} and can be retrieved later.
\end{theorem}

Here is the stored theorem again:

\getkeytheorem{stored-demo}

Finally, \pkg{keytheorems} can produce a list of theorem-like environments:

\listofkeytheorems[title={Theorems appearing in this example}]



In ordinary use, a document preamble might look like this:

\begin{verbatim}
\usepackage[shortcuts,theorems,framedtheorems]{tssuite}

\DefineScName[\Cpp]{C++}
\DefineScName{GitHub}
\DefineScNames{Python, Rust, LuaLaTeX}

\CreateTextField{\AuthorName}{author}
\CreateTextField{\ProjectTitle}{project}

\AuthorName{Alice Example}
\ProjectTitle{My Project}

\NewNumberedTheorem{theorem}{Theorem}
\NewDocumentTheorem{definition}{Definition}
\MakeFramedTheorem{ftheorem}{theorem}
\end{verbatim}

Then the body of the document can simply use the commands that were declared in
the preamble.  For example, it can mention \Cpp, \GitHub, and \Python without
repeating formatting decisions, print the project name with
\verb|\PrintTextField{project}|, and use theorem environments without repeating
setup code.

The package performs basic safety checks.  If a field is created twice, an error
is reported.  If a field is printed before it exists, an error is reported.  If
a shortcut command already exists, the package reports that instead of silently
overwriting it.  These checks are intended to make mistakes visible early,
especially in larger documents where definitions may be split across several
files.

The most useful habit is to put all package-level declarations near the top of
the preamble.  Define shortcuts first, then text fields, then theorem
environments.  Use framed theorem wrappers only for environments that need extra
visual emphasis.  For small documents, loading only the options you need keeps
the setup simple.  For lecture notes, reports, or technical manuals, using all
features together can make the source more consistent and easier to maintain.

\end{document}
